\chapter{Statistical completions and fields over APP}
\label{chap:completaciones-cuanticas}

The words Bose, Fermi, Pauli, or photon do not by themselves designate theorems of a discrete system. Obtaining them requires specification of a state space, a grading, a product, and operators. APP and TPK provide a set of oriented states and a sheet parity. Two universal completions can be constructed from those data: the symmetric algebra and the exterior algebra. This section identifies exactly what follows from the construction and separates algebraic statistics from its physical interpretation.

\section{Linear state space associated with a particle realization}

Let \(S\) be a finite set of TPK route classes at the chosen depth,
and let
\[
 \mathcal H_1=\ell^2(S)
\]
be the Hilbert space with orthonormal basis \((e_s)_{s\in S}\). A choice of
positive weights \(E_s\) defines the diagonal operator
\[
 He_s=E_se_s.
\]
The choice of \(S\), its quotient by symmetries, and the weights are part
of the model. Once these are fixed, the multiparticle completions are canonical.

\section{Symmetric and exterior completions}

We define
\[
 \mathcal F_+(\mathcal H_1)=
 \bigoplus_{n\geq0}\operatorname{Sym}^n\mathcal H_1,
 \qquad
 \mathcal F_-(\mathcal H_1)=
 \bigoplus_{n\geq0}\bigwedge^n\mathcal H_1.
\]
The first is the bosonic completion; the second, the fermionic.

\begin{theorem}[Pauli exclusion principle in the exterior completion]
For every mode \(s\in S\), the creation and annihilation operators of the
exterior completion satisfy
\[
 \{a_s,a_t^\dagger\}=\delta_{st}I,
 \qquad
 \{a_s,a_t\}=\{a_s^\dagger,a_t^\dagger\}=0.
\]
In particular, the occupation operator
\(N_s=a_s^\dagger a_s\) is a projector:
\[
 N_s^2=N_s,
\]
and its eigenvalues are \(0,1\).
\end{theorem}

\begin{proof}
Creation is the exterior product by \(e_s\); annihilation is its
adjoint. Antisymmetry implies \(e_s\wedge e_s=0\) and yields the canonical
relations by calculation on the exterior basis. Then
\[
N_s^2=a_s^\dagger a_sa_s^\dagger a_s
=a_s^\dagger(I-a_s^\dagger a_s)a_s=N_s.
\]
Every self-adjoint projector has spectrum contained in \(\{0,1\}\).
\end{proof}

The exclusion principle thus appears as a theorem of the exterior
completion. The specifically HMT question is which transport datum selects
that completion for a class of routes.

\section{Sheet grading and exchange algebra}

Suppose that the oriented sheet defines an additive character
\[
 p:\operatorname{Mor}(\mathsf P_{\rm TPK})\longrightarrow\mathbb Z/2\mathbb Z,
 \qquad p(\gamma\eta)=p(\gamma)+p(\eta).
\]
On the corresponding graded linear space, introduce the symmetry
\[
 \tau(v\otimes w)=(-1)^{p(v)p(w)}w\otimes v.
\]

\begin{theorem}[Statistical completion determined by parity]
The quotient of the tensor algebra by the relations
\[
 v\otimes w-(-1)^{p(v)p(w)}w\otimes v
\]
is symmetric on the even sector and exterior on the odd sector. Two identical
odd states have zero product; even states admit arbitrary
occupation.
\end{theorem}

\begin{proof}
If \(p(v)=p(w)=0\), the relation is \(vw=wv\). If
\(p(v)=p(w)=1\), it is \(vw=-wv\); setting \(v=w\) and working in
characteristic zero gives \(v^2=0\). Mixed sectors commute without
an additional sign. This is the universal property of the graded symmetric
algebra.
\end{proof}

The theorem converts a compositional TPK parity into an exact occupancy
rule. To obtain the physical spin--statistics theorem, one must additionally
construct a positive local Lorentzian theory and prove that sheet parity
coincides with spin under rotations. The equality does not follow solely
from both objects taking values modulo two.

\section{Sheet holonomy and \texorpdfstring{\(2\pi/4\pi\)}{2pi/4pi} periodicity}

The preceding parity admits a natural geometric realization when a
surviving extension has a closed return with a nontrivial orientation
cover. Let \(\rho\) be the return generator and let
\[
 \varepsilon:\langle\rho\rangle\longrightarrow\{-1,+1\}
\]
be its holonomy character, with \(\varepsilon(\rho)=-1\). The associated real
bundle is then the orientation line of a Möbius strip: one
turn closes the geometric projection but changes the lifted sheet; two
turns restore the orientation.

To express this property in angular terms, fix a unitary
representation \(U\) of the rotation cover on the graded space
\(\mathcal H_1=\mathcal H_{\bar0}\oplus\mathcal H_{\bar1}\), compatible with
sheet parity in the sense that
\[
 U(2\pi)v=(-1)^{p(v)}v
\]
for every homogeneous vector \(v\). Compatibility is a condition of the
realization: it identifies the generator of the orientation cover with the
character \(p\) already constructed on the routes.

\begin{theorem}[Restoration of orientation and exchange algebra]
\label{thm:dos-cuatro-pi-estadistica}
Under the preceding compatibility condition,
\[
 U(2\pi)|_{\mathcal H_{\bar0}}=I,
 \qquad
 U(2\pi)|_{\mathcal H_{\bar1}}=-I,
 \qquad
 U(4\pi)=I.
\]
Moreover, the same character \(p\) determines the graded symmetry
\[
 \tau(v\otimes w)=(-1)^{p(v)p(w)}w\otimes v.
\]
Thus the even-holonomy sector is completed by symmetric
powers and the odd-holonomy sector by exterior powers. In
the latter, the exclusion relation \(v\wedge v=0\) holds.
\end{theorem}

\begin{proof}
The first and second equalities are the compatibility condition evaluated
on each homogeneous component. Iterating one turn gives
\[
 U(4\pi)v=U(2\pi)^2v=(-1)^{2p(v)}v=v,
\]
which proves the restoration. The assertion concerning exchange is the
parity-completion theorem of the preceding section: for two even vectors
the sign is positive, whereas for two odd vectors it is negative.
\end{proof}

Three returns should be distinguished. Angular closure after \(2\pi\) is a
return of the projection; closure after \(4\pi\) restores the lifted
orientation; neither, by itself, requires every coordinate of the record
to return to its initial value. This separation is the dimensional analogue of the
HMT distinction between phase return and state return.

\subsection{The sum and product sheets as supports of the \(2\) completions}

The APP table contains two evaluations on the same support. The
additive sheet acts by translations and preserves the multiplicity of
continuations; in this sense, it accumulates. The multiplicative sheet
permutes the units and contracts fibers over noninvertible elements; in
this sense, it folds. This incompatibility precedes the statistical
choice and explains why TPK must preserve two genealogies rather than
reduce them to a single visible word.

The statistical assignment is not obtained, however, by verbally replacing
``sum'' with ``boson'' and ``product'' with ``fermion.'' The exact map is
the parity character
\[
 p:\operatorname{Mor}(\mathsf P_{\rm TPK})\longrightarrow\mathbb Z/2\mathbb Z.
\]
When the transport of an extension lies in the sector \(p=0\), its
completion admits cumulative occupation and is symmetric. When it lies in
\(p=1\), sheet reversal introduces the exchange sign and the
completion is exterior. Sum and product supply the two
local geometries; the record, orientation, and character \(p\) determine the
sector in which a route is realized.

With a diagonal Hamiltonian \(H\), energy \(\epsilon\), chemical potential
\(\mu\), and thermal parameter \(\beta\), the mean occupations of the two
sectors are then
\[
 \bar n_{\rm par}(\epsilon)
 =\frac{1}{e^{\beta(\epsilon-\mu)}-1},
 \qquad
 \bar n_{\rm impar}(\epsilon)
 =\frac{1}{e^{\beta(\epsilon-\mu)}+1}.
\]
The Bose--Einstein and Fermi--Dirac formulas are therefore the
thermodynamic evaluations of the symmetric and exterior completions
selected by transport. The periodicity \(4\pi\) and fermionic
exclusion are coordinated by the same character, but neither of them
is inferred from a mere numerical coincidence.

\section{Number and Occupation Distribution Operators in the Fock Realization}

If there are \(g\) modes and \(N\) particles, the sector dimensions are
\[
 \Omega_F(g,N)=\binom gN,
 \qquad
 \Omega_B(g,N)=\binom{g+N-1}{N}.
\]
The first formula counts subsets of modes; the second, weak compositions of
\(N\) into \(g\) parts. These counts are exact consequences of
the two completions and require no thermodynamic hypothesis.

\section{The cellular complex of the APP torus}

The APP graph is the one-skeleton of a cellular decomposition of the torus
\(T^2\) with \(9\times9\) vertices. Let
\[
 C^0\xrightarrow{d_0}C^1\xrightarrow{d_1}C^2
\]
be its real cochain complex. Choosing the uniform inner product,
define the codifferential \(\delta=d^*\) and the Laplacian
\[
 \Delta_k=d_{k-1}\delta_k+\delta_{k+1}d_k.
\]

\begin{theorem}[Harmonic sector of the digital torus]
The dimension of the space of harmonic one-forms is two:
\[
 \dim\ker\Delta_1=b_1(T^2)=2.
\]
A basis can be chosen that is represented by the uniform horizontal and
vertical flows. Both are closed and coclosed, and neither is the gradient of
a globally defined periodic function.
\end{theorem}

\begin{proof}
The Hodge theorem for finite cellular complexes identifies
\(\ker\Delta_1\) with \(H^1(T^2;\mathbb R)\). The cohomology of the torus has
rank two. Directly, the constant horizontal and vertical flows have
zero divergence and curl; their integrals over the two fundamental
cycles are independent, so they are not exact.
\end{proof}

The space of connections modulo gauge transformations is
\[
 C^1/\operatorname{im}d_0.
\]
In a finite complex equipped with the preceding inner product, each orbit
admits a Coulomb representative in \(\ker\delta_1\); within the orbit,
the choice is unique modulo the stabilizer \(\ker d_0\) of the gauge
parameter. The correct cohomological quotient is
\[
 \ker d_1/\operatorname{im}d_0
 \simeq \ker\Delta_1 .
\]
Its two harmonic classes constitute a natural mathematical candidate for
two free global modes. They are not yet the two local polarizations of
a photon in \(3+1\)-dimensional spacetime: that identification requires a
dynamical limit, a dispersion relation, and a physical gauge group.

\section{Discrete Gauss law}

For an oriented 1-cochain \(F\) and a set of vertices \(V\),
divergence is defined by
\[
 (\operatorname{div}F)(v)=
 \sum_{e:v\to w}F(e)-\sum_{e:w\to v}F(e).
\]
Summing over \(V\), every interior edge appears once with each sign and
cancels. Therefore,
\[
 \boxed{
 \sum_{v\in V}\operatorname{div}F(v)
 =\operatorname{Flux}_{\partial V}(F).}
\]
This identity is an exact Gauss law on the graph. Interpreting
\(\operatorname{div}F\) as electric charge requires the charge character and
dimensional realization constructed in Dimensional Mechanics.

\section{Reversible measurement apparatus}

Let \(S\) be a set of states and let \(r:S\to\mathbb Z/9\mathbb Z\) be a
readout. Coupling to a nonadic register is defined by
\[
 U_r:S\times\mathbb Z/9\mathbb Z\longrightarrow
 S\times\mathbb Z/9\mathbb Z,
 \qquad
 U_r(s,a)=(s,a+r(s)).
\]

\begin{proposition}[Reversibility of the coupling]
The map \(U_r\) is to permutation and
\[
 U_r^{-1}(s,a)=(s,a-r(s)).
\]
Information loss arises only after applying a reader to the
register or conditioning the state on the outcome.
\end{proposition}

\begin{proof}
The proposed inverse satisfies both composition identities by the
group law of \(\mathbb Z/9\mathbb Z\).
\end{proof}

This proposition formalizes the operational definition of measurement: coupling,
reading, and restricting compatible futures are three different maps. The
full coupling is reversible; the observational quotient need not be.

\section{Algebraic correspondences and physical realizations}

The Fock completions, exclusion conditional on parity, the harmonic
sector of APP, discrete Gauss law, and reversibility of the apparatus are complete
mathematical theorems. Their physical identifications form a diagram that
must commute:
\[
 \text{sheet parity}
 \longrightarrow\text{graded interchange}
 \longrightarrow\text{CAR/CCR}
 \longrightarrow\text{physical statistics},
\]
\[
 H^1(\mathcal A_9)
 \longrightarrow\text{discrete gauge field}
 \longrightarrow\text{massless mode}
 \longrightarrow\text{photon}.
\]
The first two arrows of each line have explicit constructions; the
last require the corresponding dynamics and physical limit. This
separation permits the bridges to be developed without conflating an analogy of
cardinality with a representation.
